\documentclass[11pt]{article}
\usepackage[margin=1in]{geometry}
\usepackage{amsmath,amssymb,amsthm,booktabs}
\usepackage[hidelinks]{hyperref}
\usepackage{microtype}

\newtheorem{theorem}{Theorem}
\newtheorem{lemma}[theorem]{Lemma}
\newtheorem{proposition}[theorem]{Proposition}
\newtheorem{corollary}[theorem]{Corollary}
\theoremstyle{definition}
\newtheorem{definition}[theorem]{Definition}
\newtheorem{remark}[theorem]{Remark}

\title{A Composite-Modulus Backbone in the Exact Type-II Erd\H{o}s--Straus Depth Spectrum}
\author{Skip Potter}
\date{Draft -- August 2026}

\begin{document}
\maketitle

\begin{abstract}
For a prime $p\equiv1\pmod4$, define the square-completion Type-II depth
\[
\delta_{\rm sq}(p)=\min\{A\ge1:\exists D\mid A^2,\ 4A-1\mid p+4D\}.
\]
The divisor-square Type-II criterion is classical; the object studied here is its exact first-hit spectrum.

We prove a symbolic composite-modulus realization theorem.  Let $q\equiv1\pmod4$ be prime and set
\[
A_q=\frac{3q+1}{4}.
\]
If $A_q$ has a prime divisor congruent to $2\pmod3$, then $4A_q-1=3q$ is composite and infinitely many primes $p\equiv1\pmod4$ satisfy
\[
\delta_{\rm sq}(p)=A_q.
\]
More generally, any divisor $D\mid A_q^2$ with $D\equiv2\pmod3$ yields such an infinite exact-depth prime family.

For every fixed $d\equiv2\pmod3$ with $3\nmid d$, primes $q$ in one reduced residue class modulo $4d$ satisfy $d\mid A_q$ and hence generate an infinite subfamily.  The choice $d=2$ gives $q\equiv5\pmod8$ and already implies
\[
N_{\rm comp}(X)\ge\left(\frac13+o(1)\right)\frac{X}{\log X}
\]
for the counting function of composite-target depths realized infinitely often.  This subfamily may be realized inside the classical hard class $p\equiv1\pmod{840}$.

We also formulate a finite CRT certificate theorem for arbitrary composite target layers and record computationally certified full-layer examples at depths $649$ and $22339$.  The results concern Type-II exact depth only and do not prove universal Type-II solvability or the Erd\H{o}s--Straus conjecture.
\end{abstract}

\section{Introduction}
The Erd\H{o}s--Straus conjecture asks whether
\[
\frac4n=\frac1x+\frac1y+\frac1z
\]
has a solution in positive integers for every $n\ge2$.  For prime denominators, standard solutions split into Type I and Type II.  We restrict attention to Type II.

A classical divisor-square coordinate for prime Type II asks for positive integers $A,D$ satisfying
\[
D\mid A^2,\qquad 4A-1\mid p+4D. \tag{1}
\]
This mechanism belongs to the historical Th\'epault--Mizony line and is not claimed here.  Modern complete Type-II criteria appear, for example, in the work of Elsholtz--Tao and Pomerance--Weingartner; recent work of Dahan studies Type-II search depth.

For primes $p\equiv1\pmod4$, define
\[
\delta_{\rm sq}(p)
=\min\{A\ge1:\exists D\mid A^2\text{ satisfying (1)}\}, \tag{2}
\]
with value $\infty$ if there is no Type-II witness.  Recent public work of FCF/CENTL studies the resulting exact-depth spectrum, proving an infinite prime-modulus backbone and identifying the composite-modulus spectrum as a dependency-rich residual problem.  Our contribution is an elementary infinite backbone lying entirely in that composite target regime.

The proof uses one universal safe residue and the fact that $3$ is the smallest positive integer congruent to $3\pmod4$.  These two observations force complete avoidance of every earlier Type-II layer while retaining a prescribed target witness.

\section{The square-completion coordinate}
We first record the standard completeness statement.

\begin{proposition}[complete square-divisor coordinate]
For a prime $p\equiv1\pmod4$, a standard Type-II Erd\H{o}s--Straus solution exists if and only if there are positive integers $A,D$ such that
\[
D\mid A^2,\qquad p\nmid D,\qquad 4A-1\mid p+4D.
\]
\end{proposition}

\begin{remark}
This proposition is prior mathematics.  It can be obtained from the classical Type-II parametrizations or from the divisor-square completion of the Th\'epault condition.  We use it only to justify that $\delta_{\rm sq}$ measures the complete standard prime Type-II family rather than a sufficient subfamily.
\end{remark}

\section{A universal safe residue}
The following elementary lemma is the engine of the construction.

\begin{lemma}[universal safe residue]\label{lem:safe}
For every $B\ge1$ and every $E\mid B^2$,
\[
4B-1\nmid4E+1.
\]
Equivalently, the residue $1$ modulo $4B-1$ is never a square-completed Type-II trap residue at depth $B$.
\end{lemma}

\begin{proof}
Let
\[
g=\gcd(B,E),\qquad B=ga,\qquad E=gd,\qquad\gcd(a,d)=1.
\]
Because $E\mid B^2$, we have $d\mid g$; write $g=dh$.  Thus
\[
B=dha,\qquad E=d^2h.
\]
Suppose $4B-1\mid4E+1$.  Modulo $4B-1$ we have $4B\equiv1$ and $4E\equiv-1$, hence
\[
E\equiv-B\pmod{4B-1}.
\]
Therefore
\[
4B-1\mid E+B=dh(d+a).
\]
Every prime divisor of $dh$ divides $B$, so $\gcd(4B-1,dh)=1$.  Hence
\[
4B-1\mid d+a.
\]
But
\[
4B-1=4dha-1>d+a,
\]
a contradiction.
\end{proof}

\section{The composite-modulus backbone}
We now state the main theorem.

\begin{theorem}[structural composite exact-depth criterion]\label{thm:main}
Let $q\equiv1\pmod4$ be prime and put
\[
A_q=\frac{3q+1}{4}.
\]
Suppose there is a divisor
\[
D\mid A_q^2,\qquad D\equiv2\pmod3.
\]
Then
\[
4A_q-1=3q
\]
is composite and there are infinitely many primes $p\equiv1\pmod4$ such that
\[
\boxed{\delta_{\rm sq}(p)=A_q.}
\]
The depth-$A_q$ witness is the chosen divisor $D$.
\end{theorem}

\begin{proof}
Set
\[
m=3q=4A_q-1,
\qquad a\equiv-4D\pmod m.
\]
Since $D\equiv2\pmod3$,
\[
a\equiv1\pmod3. \tag{3}
\]
Every integer $n\equiv a\pmod m$ has a depth-$A_q$ witness.

Let $B<A_q$ and write $m_B=4B-1$.  We claim $q\nmid m_B$.  Otherwise $m_B=qs$.  Since $m_B\equiv3\pmod4$ and $q\equiv1\pmod4$, one has $s\equiv3\pmod4$ and hence $s\ge3$.  Thus
\[
m_B\ge3q=m,
\]
contradicting $B<A_q$.  Therefore
\[
\gcd(m,m_B)\mid3 \tag{4}
\]
for every earlier depth.

Let
\[
L=\operatorname{lcm}_{1\le B<A_q}(4B-1).
\]
By (3)--(4), the system
\[
n\equiv-4D\pmod{3q},
\qquad
n\equiv1\pmod L,
\qquad
n\equiv1\pmod4 \tag{5}
\]
is CRT-compatible.  Let its solution be $n\equiv n_0\pmod M$, where
\[
M=\operatorname{lcm}(4,3q,L).
\]

For every $B<A_q$, (5) gives $n\equiv1\pmod{4B-1}$.  Lemma~\ref{lem:safe} therefore excludes every earlier square-completed Type-II witness.  At the target depth, $n\equiv-4D\pmod{3q}$ and $D\mid A_q^2$, so a witness exists.

It remains to check reducedness.  Every prime divisor contributed by $L$ sees residue $1$.  The only new odd prime factor of the target modulus is $q$.  Because $q\nmid A_q$, no divisor of $A_q^2$ is divisible by $q$; hence $q\nmid D$ and $-4D\not\equiv0\pmod q$.  The residue modulo $4$ is $1$.  Thus $\gcd(n_0,M)=1$.

Dirichlet's theorem gives infinitely many primes in (5).  Every such prime is congruent to $1\pmod4$, has no Type-II witness below $A_q$, and has the prescribed witness at $A_q$.
\end{proof}

\begin{corollary}[prime-factor form]\label{cor:factor}
If $q\equiv1\pmod4$ is prime and $A_q=(3q+1)/4$ has a prime divisor
\[
\ell\equiv2\pmod3,
\]
then $A_q$ is realized as $\delta_{\rm sq}(p)$ for infinitely many primes $p$.
\end{corollary}

\begin{proof}
Take $D=\ell$ in Theorem~\ref{thm:main}.
\end{proof}

\begin{remark}
If every prime divisor of $A_q$ is $1\pmod3$, Theorem~\ref{thm:main} does not apply.  This is an obstruction only to the present universal-safe-residue construction; it is not asserted to be an obstruction to exact-depth realization by other mechanisms.
\end{remark}

\section{Infinite fixed-witness subfamilies}
The structural theorem immediately contains infinitely many explicit arithmetic families.

\begin{corollary}[fixed witness]\label{cor:fixed}
Fix $d\ge1$ with
\[
d\equiv2\pmod3,
\qquad 3\nmid d.
\]
There is a unique reduced residue class $a_d\pmod{4d}$ satisfying
\[
3a_d\equiv-1\pmod{4d}.
\]
For every prime
\[
q\equiv a_d\pmod{4d},
\]
the depth
\[
A_q=\frac{3q+1}{4}
\]
is realized exactly by infinitely many primes, with witness $D=d$.
\end{corollary}

\begin{proof}
The congruence implies $d\mid A_q$.  Thus $d\mid A_q^2$, while $d\equiv2\pmod3$.  Apply Theorem~\ref{thm:main}.  The class $a_d$ is reduced because $3$ is invertible modulo $4d$.
\end{proof}

Dirichlet's theorem gives infinitely many primes $q$ in each fixed-witness family.

\subsection{The $D=2$ backbone}
Taking $d=2$ gives
\[
q\equiv5\pmod8.
\]
Thus
\[
\left\{\frac{3q+1}{4}:q\text{ prime},\ q\equiv5\pmod8\right\}
\subseteq\mathcal D^{\rm comp}_{\infty}. \tag{6}
\]
This single subfamily already has the correct prime-order scale.

\begin{corollary}[quantitative lower bound]\label{cor:count}
Let
\[
N_{\rm comp}(X)=\#\{A\le X:4A-1\text{ composite and }\delta_{\rm sq}(p)=A\text{ for infinitely many primes }p\}.
\]
Then
\[
\boxed{
N_{\rm comp}(X)\ge
\left(\frac13+o(1)\right)\frac{X}{\log X}.
}
\]
\end{corollary}

\begin{proof}
By (6), it suffices to count primes $q\equiv5\pmod8$ with $q\le(4X-1)/3$.  The prime number theorem in arithmetic progressions gives
\[
\pi(Y;8,5)\sim\frac14\operatorname{Li}(Y),
\]
and $Y\sim4X/3$.
\end{proof}

\begin{corollary}[hard-core realization]
For every prime $q\equiv5\pmod8$ with $q>5$, infinitely many primes
\[
p\equiv1\pmod{840}
\]
have exact square-completion depth $(3q+1)/4$.
\end{corollary}

\begin{proof}
For $q>5$, $\gcd(3q,840)=3$, while the target residue $-8$ and the hard residue $1$ agree modulo $3$.  Thus the congruence $n\equiv1\pmod4$ in (5) may be strengthened to $n\equiv1\pmod{840}$.
\end{proof}

\section{Finite certificates for arbitrary composite layers}
The preceding theorem uses the special target modulus $3q$.  For a general composite target modulus, exact first-hit can be reduced to a finite congruence certificate.

Fix $A,D$ with $D\mid A^2$, set $m_A=4A-1$ and $a\equiv-4D\pmod{m_A}$, and write
\[
n=a+m_Ak. \tag{7}
\]
An earlier witness $(B,E)$, $B<A$, $E\mid B^2$, imposes
\[
n\equiv-4E\pmod{m_B},
\qquad m_B=4B-1. \tag{8}
\]
Put $g=\gcd(m_A,m_B)$.  The two progressions intersect exactly when
\[
a\equiv-4E\pmod g.
\]
When they do, (7)--(8) remove one class
\[
k\equiv c_{B,E}\pmod{m_B/g}. \tag{9}
\]
Thus all earlier layers give a finite family of forbidden congruence classes in the single parameter $k$.

\begin{theorem}[finite exact-depth certificate]
Suppose the forbidden system (9) admits a residue class
\[
k\equiv r\pmod L
\]
which avoids every forbidden congruence and for which
\[
\gcd(a+m_Ar,m_AL)=1.
\]
Then infinitely many primes have exact square-completion depth $A$ and target residue $a\pmod{m_A}$.
\end{theorem}

\begin{proof}
Every integer in the reduced progression
\[
n\equiv a+m_Ar\pmod{m_AL}
\]
has the target witness and avoids every earlier witness.  Apply Dirichlet's theorem.
\end{proof}

In practice one need not print the potentially enormous CRT integer $r$.  The modulus $L$ can be factored into maximal prime-power coordinates and a compact assignment can be checked coordinate by coordinate.

\section{Two certified full composite layers}
The repository accompanying this work contains deterministic verifiers for two larger composite target layers.

At
\[
A=649=11\cdot59,
\qquad
4A-1=2595=3\cdot5\cdot173,
\]
the realized divisor-modulus primitive branches are represented by
\[
D\in\{11,59,7139,38291\}.
\]
Each branch admits a reduced finite exact-depth certificate, so every one supports infinitely many primes with $\delta_{\rm sq}(p)=649$.

At
\[
A=22339=89\cdot251,
\qquad
4A-1=89355=3\cdot5\cdot7\cdot23\cdot37,
\]
the four realized primitive branches
\[
D\in\{89,251,22361,67079\}
\]
likewise admit reduced exact-depth certificates.  Hence every one supports infinitely many primes with $\delta_{\rm sq}(p)=22339$.

These examples are not needed for the symbolic backbone theorem.  They show that the composite exact-depth spectrum extends beyond the elementary $3q$ construction and can exhibit full-layer realization in more dependency-rich moduli.

\section{Relation to prior work and claim boundary}
The divisor-square Type-II mechanism itself is prior art.  Historical provenance reaches at least to Th\'epault, with later presentations by Mizony and Gardes.  Modern complete Type-II parametrizations include those of Elsholtz--Tao and Pomerance--Weingartner.  Dahan studies Type-II sieve dimension and search depth.

Recent public FCF/CENTL work develops an exact-depth spectrum, a prime-modulus backbone, structural shadow theory, and a composite residual core for the square-completed Type-II route.  The present theorem is intended as a contribution inside that composite residual spectrum: it supplies an explicit symbolic family with target modulus $3q$, a structural factor criterion, infinitely many fixed-witness arithmetic subfamilies, and the lower bound in Corollary~\ref{cor:count}.

Targeted literature searches through August 30, 2026 did not locate this exact family or structural criterion.  This negative search does not establish priority.  A broader historical review, including non-digitized and non-arXiv literature, remains necessary before any claim of first discovery.

Nothing proved here implies universal Type-II solvability.  In particular, the Erd\H{o}s--Straus conjecture remains open.

\section{Open questions}
The construction suggests several natural problems.
\begin{enumerate}
\item Determine the density of primes $q\equiv1\pmod4$ for which $A_q=(3q+1)/4$ has a prime factor $2\pmod3$.  A sieve showing that the complementary split-prime-only values are sparse would sharpen Corollary~\ref{cor:count}.
\item Classify further cofactors replacing $3$.  The proof shows why $3$ is exceptional: for $r>3$, a prime factor $q$ of a target modulus $rq$ can already appear in the earlier modulus $3q$, so a new compatibility mechanism is needed.
\item Determine whether there are infinitely many composite layers for which every realized divisor-modulus primitive branch is globally first-hit, as happens at $649$ and $22339$.
\item Compare the square-completion depth with Dahan's pair depth and with the minimal López Type A/B depth spectrum.
\end{enumerate}

\begin{thebibliography}{9}
\bibitem{ET}
C. Elsholtz and T. Tao,
Counting the number of solutions to the Erd\H{o}s--Straus equation on unit fractions,
\emph{Journal of the Australian Mathematical Society} 94 (2013), 50--105.

\bibitem{PW}
C. Pomerance and A. Weingartner,
Exceptions to the Erd\H{o}s--Straus--Schinzel conjecture,
\emph{Ramanujan Journal} 69 (2026).

\bibitem{Dahan}
B. Dahan,
Sieve dimension and search depth for the Erd\H{o}s--Straus conjecture, $n\equiv1\pmod{24}$,
arXiv:2608.24035, 2026.

\bibitem{MG}
L. Mizony and J.-L. Gardes,
Un point sur la conjecture d'Erd\H{o}s et Straus,
Universit\'e Claude-Bernard Lyon 1, 2012.

\bibitem{FCF}
Free Computation Foundation / CENTL,
Erd\H{o}s--Straus research library: square-completion prior art, exact-depth spectrum, and current frontier,
2026.
\end{thebibliography}

\end{document}
